\ncdlesson{13}{H.C. Andersens skygge: Eventyr for børn eller satire for voksne?}{安徒生的影子：儿童童话还是成人讽刺？}
{DU3 Modul 5 (PD3) $\rightarrow$ Modul 6 (Studieprøven, B2 $\rightarrow$ C1)｜对标 NCE 4 \& \textit{Litterær Kritik, Satire \& Undertekst}}
{\textbf{文学批评与文本潜台词分析语体}（双关多义 \textit{tvetydighed} · 文本双层意涵 \textit{dobbeltbund} · 潜台词 \textit{undertekst} · 阶级讽刺 \textit{klassesatire} · 隐喻寓言 \textit{allegori}）· 文学评论核心动词搭配（\textit{afspejle, udstille, gennemskue, fortolke som}）与表里对比句式}
{离开丹麦后，人们提起汉斯·克里斯蒂安·安徒生（H.C. Andersen），脑海中浮现的往往是温馨的睡前童话或动画改编。然而在丹麦文学批评（\textit{litteraturkritik}）与 Studieprøven 文化分析写作中，读懂安徒生恰恰意味着揭开“天真表层”之下的\textbf{文本双层意涵（\textit{dobbeltbund}）与成人讽刺}。本课带你精读《皇帝的新衣》《影子》与《丑小鸭》的深层潜台词，掌握客观、敏锐的丹麦语文学评论语体。}

% ---------------------------------------------------------------------
\begin{foerduleaeser}
带着下面两个问题阅读 Tekst A 和 Tekst B：
\begin{enumerate}
  \item \textit{Hvorfor mener litteraturforskere i Tekst A, at H.C. Andersens eventyr rummer en ”dobbeltbund”, og hvad udstiller både ”Kejserens nye klæder” og ”Skyggen” under overfladen?}（根据 Tekst A，为什么文学学者认为安徒生的童话包含“双层意涵”？《皇帝的新衣》和《影子》在表层之下分别揭露了什么？）
  \item \textit{Hvilken uenighed opstår mellem Henrik og Sofie i Tekst B, når de diskuterer, hvordan ”Den grimme ælling” bør læses?}（在 Tekst B 中，当亨里克与索菲娅讨论应如何解读《丑小鸭》时，两人产生了什么分歧？）
\end{enumerate}
\end{foerduleaeser}

% ---------------------------------------------------------------------
\ncdsubsec{二、Tekster（文学批评专论与播客对谈双语篇）}

\subsection*{Tekst A：\textit{Bag den uskyldige overflade: Den skarpe satire hos H.C. Andersen}（文学评论专论 · 约 360 词）}

\begin{tipbox}{语篇说明}
本篇是一篇标准的 C1 文学评论短论，剖析安徒生童话如何同时面向儿童与成人读者。请重点关注\textbf{\color{NCDNavy}深蓝色标记的文学批评术语与分析动词}（\textit{dobbeltbund, undertekst, klassesatire, allegori, afspejle, udstille, gennemskue}），以及\textbf{\color{NCDRed}深红色标记的表层与深层对照句式}。
\end{tipbox}

Uden for Danmark bliver H.C. Andersen ofte opfattet som en harmløs fortæller af søde godnathistorier for børn, hvor alting til sidst falder trygt på plads. Blandt danske litteraturforskere hersker der imidlertid bred enighed om, at hans eventyr er skrevet med en bevidst \da{dobbeltbund}. Allerede i sin egen levetid understregede Andersen, at han fortalte for både børn og voksne, og \dabold{under den tilsyneladende uskyldige overflade gemmer der sig} en skarp \da{tvetydighed}, der kræver, at læseren evner at læse mellem linjerne.

Tag for eksempel \textit{Kejserens nye klæder} fra 1837. Ved første øjekast ligner fortællingen en humoristisk fabel om en narragtig hersker, der lader sig narre af to snedige vævere. \dabold{Ved nærmere eftersyn fungerer eventyret snarere som} en bidende \da{klassesatire} over konformitet, smiger og intellektuel fejhed i magtens korridorer. Ingen ved hoffet tør indrømme, at de ikke kan se det usynlige stof, af frygt for at blive stemplet som dumme eller uegnede til deres høje embede. Dermed \da{udstiller} teksten et kollektivt \da{selvbedrag}, hvor kejserens \da{forfængelighed} og omgivelsernes tavshed gensidigt forstærker hinanden, indtil et barn bryder illusionen med en enkel, ufiltreret sandhed.

Endnu mørkere bliver denne \da{undertekst} i \textit{Skyggen} fra 1847, der ofte regnes for et af Andersens mest \da{uhyggelige} og \da{flertydige} værker. Her mister en godtroende lærd sin egen skygge under en rejse i syden, hvorefter skyggen år senere vender tilbage som en velhavende og kynisk verdensmand. Skridt for skridt bytter de to skikkelser roller: den tomme form overtager magten, gifter sig med en prinsesse og lader til sidst sin tidligere herre henrette for at skjule sandheden om sin oprindelse. \dabold{Snarere end blot at være} en fantastisk spøgelseshistorie læses \textit{Skyggen} i dag som en eksistentiel \da{allegori} over modernitetens overfladiskhed, hvor kynisme og opportunistisk imagepleje kan fortrænge ægte integritet.

Selv et elsket værk som \textit{Den grimme ælling} \da{afspejler} en social smerte, som rækker langt ud over børneværelset. Fortællingen kan ganske vist læses som en trøstende historie om at finde sit sande jeg, men den rummer samtidig en bitter kritik af 1800-tallets stive klasseskel og snæversynede andegård. Andersens kunstneriske storhed ligger netop i denne uforløste spænding mellem barnlig naivitet og voksen \da{spidsfindighed}: den, der kun ser det trygge eventyr, overser halvdelen af værket.

\begin{transbox}
\small
\textbf{\color{NCDTeal}[Tekst A 中文精译]}\\
在丹麦以外的地方，安徒生常常被看作一位人畜无害、专门为儿童讲述甜蜜睡前故事的作家——在他的故事里，一切最终都会安稳妥帖地各归其位。然而在丹麦文学学者之间却存在广泛共识：他的童话是带着自觉的“双层意涵”写就的。早在生前，安徒生就曾强调自己既为孩子讲故事，也为成年人写作；而在看似天真无邪的表层之下，实则隐藏着敏锐的双关多义，要求读者具备“读出字里行间弦外之音”的能力。

以 1837 年的《皇帝的新衣》为例。乍看之下，这个故事像是一则幽默寓言，讲述一位愚蠢的统治者如何被两个狡猾的织工欺骗。但经仔细审视，这篇童话更像是一部辛辣的阶级讽刺之作，直指权力走廊中的从众心理、阿谀奉承与智识怯懦。宫廷里没有人敢承认自己看不见那块隐形布料，因为人人都害怕被贴上愚蠢或不称职于自己高位的标签。文本由此暴露出一种集体自我欺骗：皇帝的虚荣心与周围人的沉默彼此强化，直到一个孩子用一句简单而未经修饰的真话打破幻象。

这种潜台词在 1847 年的《影子》中变得更加幽暗，该作常被视为安徒生最令人毛骨悚然且意涵繁复的作品之一。故事里，一位轻信他人的学者在南方旅行时失去了自己的影子，多年后影子却以一位富有而犬儒的世故名流身份归来。一步步地，这两个形象互换了角色：空洞的躯壳篡夺了权力，迎娶了公主，最终甚至下令处决昔日的主人，以掩盖自己出身的真相。《影子》与其说仅仅是一个奇幻鬼故事，不如说在今天被解读为一则关于现代性浅薄本质的存在主义寓言——在这个世界里，犬儒主义与投机取巧的形象包装能够排挤并取代真正的品格操守。

就连《丑小鸭》这样广受喜爱的作品，也折射出一种远远超出儿童房范畴的社会痛感。这个故事固然可以被读作寻找真实自我的心灵慰藉，但与此同时，它也包含着对十九世纪僵化阶级壁垒以及狭隘“鸭群世界”的苦涩批判。安徒生的艺术伟大之处，恰恰在于这种孩童般的纯真与成人般的精微敏锐之间那未曾消解的张力：只看得见令人安心的童话表层的人，便错过了这部作品的一半。
\end{transbox}

\subsection*{Tekst B：\textit{Debat i litteraturpodcasten: Bør ”Den grimme ælling” læses som trøst eller samfundskritik?}（文学播客对谈 · 约 140 词）}

\textbf{阅读任务（理解与核对）：}\textit{Læs samtalen mellem litteraturformidleren Henrik og gymnasielæreren Sofie, og besvar spørgsmålene nedenfor.}

\begin{emailbox}{播客对谈｜Henrik（文学节目主持）\& Sofie（高中丹麦语教师）}
\textbf{Henrik:} Sofie, mange voksne husker stadig \textit{Den grimme ælling} som en varm fortælling om håb og tålmodighed. Men når vi læser den igen i dag, virker den næsten lidt ubarmhjertig, ikke?\\
\textbf{Sofie:} Jo, det gør den faktisk. \dabold{Hvor} mange opfatter historien som en simpel trøst, \dabold{rummer} teksten jo en tydelig \da{klassesatire}. Ællingen bliver slet ikke accepteret for sin karakter eller venlighed, men først respekteret, da de andre opdager, at den tilhører en finere art.\\
\textbf{Henrik:} Præcis --- hele moralen hviler næsten på, at man er født som svane fra begyndelsen. Betyder det så, at eventyret i virkeligheden er elitært?\\
\textbf{Sofie:} Det kan godt virke sådan ved første øjekast. Alligevel bør man \da{fortolke} eventyret \da{som} både et selvforsvar fra en fattig skomagersøn og en skarp kritik af snobberiet i datidens København. Hvis vi \da{gennemskuer} den \da{tvetydighed}, bliver Andersens \da{undertekst} langt mere spændende end nogen enkel moral.
\end{emailbox}

\textbf{Tekst B 阅读核对（Rigtigt / Forkert / Kort svar，答案见第六部分）：}
\begin{enumerate}
  \item \textit{Henrik og Sofie er enige om, at ”Den grimme ælling” udelukkende er en varm og simpel børnehistorie.} \underline{\hspace{2.2cm}}
  \item \textit{Hvornår bliver ællingen ifølge Sofie endelig respekteret af de andre i andegården?} \underline{\hspace{2.2cm}}
  \item \textit{Hvilke to ting kan man ifølge Sofie fortolke eventyret som?} \underline{\hspace{2.2cm}}
\end{enumerate}

\begin{tipbox}{本课发音与语调要点（Udtale \& Prosodi）}
\begin{itemize}
  \item \textbf{文学批评核心词的重音与静音字母}：\textit{en tvetydighed} \ipa{tˢveˈtˢyːðiˌheðˀ}（重音在第二音节 \textit{-ty-}，词尾带弱化 \textit{stød}）；\textit{tilsyneladende} \ipa{tˢelˈsyːnəˌlæːˀðənə}；\textit{en forfængelighed} \ipa{fʌˈfɛŋˀəliˌheðˀ}（\textit{-lig-} 中的 \textit{g} 不发音）；\textit{uhyggelig} \ipa{uˈhyɡəli}（词尾 \textit{g} 不发音！）。
  \item \textbf{外来文学学术词重音后移}：\textit{en allegori} \ipa{aleɡoˈʁiːˀ}（重音与 \textit{stød} 落在最末长元音 \textit{-ri} 上！）；\textit{en satire} \ipa{saˈtˢiːɐ}（重音在 \textit{-ti-}）；\textit{spidsfindig} \ipa{sb̥esˈfenˀi}（\textit{spids} 中的 \textit{d} 不发音，词尾 \textit{g} 不发音）。
  \item \textbf{对照句型的语流重音}：朗读 \textit{Ved første øjekast... Ved nærmere eftersyn...} 时，对比重音必须清晰落在 \textit{\textbf{før}ste} 与 \textit{\textbf{nær}mere} 上，借音调反差体现“表层误解 vs 深度洞见”的批评转折。
\end{itemize}
\end{tipbox}

\begin{transbox}
\small
\textbf{\color{NCDTeal}[Tekst B 中文对照]}\\
\textbf{亨里克：}索菲娅，许多成年人至今仍把《丑小鸭》记作一个关于希望与耐心的温暖故事。但当我们今天重读它时，它给人的感觉几乎有点冷酷无情，不是吗？\\
\textbf{索菲娅：}是的，确实如此。人们往往把这则故事看作简单的安慰，但文本中其实包含着明显的阶级讽刺。小鸭子根本不是因为自己的品格或善良而被接纳，而是直到其他动物发现它本就属于更高贵的物种时，才开始受到尊重。\\
\textbf{亨里克：}没错——整篇故事的寓意几乎建立在“你从一开始就天生是天鹅”这一点上。那这是不是意味着，这篇童话实际上带有精英主义色彩？\\
\textbf{索菲娅：}乍看之下确实可能给人这种印象。尽管如此，我们还是应当把这篇童话既解读为一位贫苦鞋匠之子的自我辩护，也解读为对当时哥本哈根势利风气的尖锐批判。只要我们看穿这种双关多义，安徒生的潜台词就会比任何单一的道德说教都精彩得多。
\end{transbox}

% ---------------------------------------------------------------------
\ncdsubsec{三、Gloser（核心生词表）}

\begin{tipbox}{词表重点}
文学分析极其讲究“动作指向与搭配”：表示“把某物解读为……”固定使用 \textit{at fortolke noget \dabold{som} noget}（而非 \textit{til}）；表示“对某现象的寓言”用 \textit{en allegori \dabold{over / på} noget}；表示“揭露/曝露虚伪”则使用动词 \textit{at udstille} 或 \textit{at gennemskue}。此外，\textit{tvetydighed} 与 \textit{flertydig} 中词根 \textit{tyd-} 均来自 \textit{betyde}（意指），指“多重含义”，切勿与表示“怀疑”的 \textit{tvivl} 混淆。
\end{tipbox}

\begin{glosertable}
名词 & en skygge & \textit{(skyggen, skygger, skyggerne)} & 影子，阴影 \\
名词 & et eventyr & \textit{(eventyret, eventyr, eventyrene)} & 童话，奇遇故事 \\
名词 & en tvetydighed & \textit{(tvetydigheden, tvetydigheder, tvetydighederne)} & 双关多义，模棱两可 \\
名词 & en undertekst & \textit{(underteksten, undertekster, underteksterne)} & 潜台词，言外之意 \\
名词 & en dobbeltbund & \textit{(dobbeltbunden, kun ental)} & 文本双层意涵，深层底蕴 \\
名词 & en allegori & \textit{(allegorien, allegorier, allegorierne)} & 寓言，隐喻象征作品 \\
名词 & en klassesatire & \textit{(klassesatiren, klassesatirer, klassesatirerne)} & 阶级讽刺（作品） \\
名词 & en forfængelighed & \textit{(forfængeligheden, kun ental)} & 虚荣，爱慕虚荣 \\
名词 & et selvbedrag & \textit{(selvbedraget, selvbedrag, selvbedragene)} & 自我欺骗，自欺 \\
名词 & en ælling & \textit{(ællingen, ællinger, ællingerne)} & 小鸭，雏鸭 \\
名词 & en spidsfindighed & \textit{(spidsfindigheden, spidsfindigheder, spidsfindighederne)} & 精微机敏，尖锐讽刺感 \\
\midrule
动词 [har] & at afspejle & \textit{(afspejler, afspejlede, har afspejlet)} & 折射，反映出 \\
动词 [har] & at udstille & \textit{(udstiller, udstillede, har udstillet)} & 暴露，揭露（缺点/虚伪） \\
动词 [har] & at fortolke & \textit{(fortolker, fortolkede, har fortolket)} & 解读，阐释 \\
动词 [har] & at gennemskue & \textit{(gennemskuer, gennemskuede, har gennemskuet)} & 看穿，识破 \\
\midrule
形容词 & tilsyneladende & \textit{(tilsyneladende, 不变)} & 表面上的，看似的 \\
形容词 & flertydig & \textit{(flertydigt, flertydige)} & 多义的，含意丰富的 \\
形容词 & uhyggelig & \textit{(uhyggeligt, uhyggelige)} & 阴森的，令人不寒而栗的 \\
\end{glosertable}

% ---------------------------------------------------------------------
\ncdsubsec{四、Grammatik \& Stilistik（文学批评修辞与潜台词分析精讲）}

\begin{grammatikbox}{修辞支柱 1：由表及里的文学批评“双层对照”句型框架}
写 C1 级别文学或文化评论时，最核心的推进逻辑是\textbf{先交代大众对文本的“表层误读”，再翻转切入“深层潜台词分析”}。请熟记以下四组高阶句型公式：

\vspace{0.3em}
\renewcommand{\arraystretch}{1.24}
{\small
\begin{tabularx}{\linewidth}{@{} >{\bfseries\raggedright\arraybackslash}p{3.0cm} >{\raggedright\arraybackslash}p{5.2cm} >{\raggedright\arraybackslash}X @{}}
\toprule
\textbf{分析维度} & \textbf{核心句型公式（注意 V2 倒装）} & \textbf{Tekst A/B 示范句} \\
\midrule
\da{1. 表层印象 vs 深度洞见} &
\textit{\dabold{Ved første øjekast} ligner X...}\newline
\textit{\dabold{Ved nærmere eftersyn} fungerer X snarere som...} &
\textit{Ved nærmere eftersyn fungerer eventyret snarere som en bidende klassesatire.} \\
\midrule
\da{2. 否定单一读法并升级} &
\textit{\dabold{Snarere end blot at} være A, læses værket som B.} &
\textit{Snarere end blot at være en spøgelseshistorie læses Skyggen som en allegori.} \\
\midrule
\da{3. 揭示潜台词与张力} &
\textit{Under den tilsyneladende... gemmer der sig...}\newline
\textit{Teksten rummer en bevidst dobbeltbund.} &
\textit{Under den uskyldige overflade gemmer der sig en skarp tvetydighed.} \\
\midrule
\da{4. 对比大众与专业视角} &
\textit{\dabold{Hvor} mange opfatter X som..., \dabold{peger} forskere på...} &
\textit{Hvor mange opfatter historien som en trøst, rummer teksten jo en klassesatire.} \\
\bottomrule
\end{tabularx}
}
\end{grammatikbox}

\begin{grammatikbox}{语法点 2：文学批评的核心句法结构与长句凝练技巧}
\begin{enumerate}
  \item \textbf{介词短语置于句首时的 V2 倒装铁律}：\\
    \textit{Ved første øjekast}, \textit{Ved nærmere eftersyn}, \textit{Under overfladen} 以及 \textit{Snarere end blot at + 不定式} 放在句首充当状语成分时，后面主句的第一位必须是\textbf{变位动词}（丹麦语不在这些句首状语后打逗号）：\\
    \ding{51}\ \textit{Under overfladen \dabold{gemmer} der sig en skarp kritik.}（\ding{55}\ \textit{*Under overfladen, der gemmer sig...}）\\
    \ding{51}\ \textit{Snarere end blot at underholde \dabold{udfordrer} eventyret læserens syn på magt.}
  \item \textbf{客观引述与被动评价句型（避免滥用“我认为”）}：\\
    文学学术文章常借助非人称或被动表达提升客观权威感：\\
    \ding{172} \textit{Værket \da{regnes ofte for} at være...}（该作常被视为……）\\
    \ding{173} \textit{Fortællingen \da{kan læses som} en allegori over...}（故事可被解读为关于……的寓言）\\
    \ding{174} \textit{Blandt litteraturforskere \da{hersker der enighed om, at}...}（学者之间普遍认同……）
  \item \textbf{形式主语 \textit{der} + 存现/隐藏动词}：表达“字里行间隐藏着某种深意”时，用 \textit{gemmer der sig + 不定名词} 或 \textit{ligger der + 不定名词}，将信息焦点自然留到句末。
\end{enumerate}
\end{grammatikbox}

\begin{grammatikbox}{文学批评避坑清单（Top 4 Pitfalls）}
\begin{enumerate}
  \item \textbf{\ding{172} \textit{fortolke} 与 \textit{læse} 表“解读为”时只能接 \dabold{som}}：\ding{55}\ \textit{*Man kan fortolke Skyggen til en kritik af magt.} $\rightarrow$ \ding{51}\ \textit{Man kan fortolke Skyggen \da{som} en kritik af magt.}
  \item \textbf{\ding{173} \textit{en allegori} 后接 \dabold{over} 或 \dabold{på}，切勿直译英文 \textit{of}}：\ding{51}\ \textit{en allegori \da{over / på} modernitetens overfladiskhed}（而非 \ding{55}\ \textit{*en allegori af}）。
  \item \textbf{\ding{174} 区分 \textit{en skæbne}（命运）与 \textit{et eventyr}（童话/奇遇）的词性}：\textit{eventyr} 是中性名词，单复同形（\textit{et eventyr, alle hans eventyr}），切勿受“很多童话”影响误加复数词尾 \ding{55}\ \textit{*eventyrer}（\textit{en eventyrer} 意思是“冒险家”！）。
  \item \textbf{\ding{175} 区分 \textit{snarere end} 与 \textit{hellere end}}：\textit{hellere end} 表达个人偏好（宁愿）；文学批评中表达“与其说是 A，不如说是 B”的客观重新定性，必须用 \da{snarere end}。
\end{enumerate}
\end{grammatikbox}

% ---------------------------------------------------------------------
\ncdsubsec{五、Mønsterøvelser（文学分析句型操练）}

\textbf{练习 A：文学批评核心词汇与介词搭配精准填空（每空填入最准确的 1 个丹麦语词）}
\begin{enumerate}
  \item Under den uskyldige overflade rummer Andersens eventyr en bevidst \underline{\hspace{2.2cm}}.
  \item Ved nærmere eftersyn fungerer \textit{Kejserens nye klæder} som en bidende \underline{\hspace{2.2cm}} over magtens smiger.
  \item Teksten \underline{\hspace{2.2cm}} et kollektivt selvbedrag, hvor ingen ved hoffet tør sige sandheden højt.
  \item Man kan fortolke \textit{Skyggen} \underline{\hspace{2.2cm}} en eksistentiel advarsel mod tom opportunisme.
  \item Eventyret \underline{\hspace{2.2cm}} på én gang tidens stive klasseskel og forfatterens egen erfaring med social udstødelse.
  \item Snarere \underline{\hspace{2.2cm}} blot at trøste børn stiller fortællingen ubehagelige spørgsmål til de voksne.
  \item Hvis man kun læser handlingen bogstaveligt, overser man let den kritiske \underline{\hspace{2.2cm}} mellem linjerne.
  \item Først når læseren \underline{\hspace{2.2cm}} ironien i slutningen, åbner værket sig i hele sin dybde.
\end{enumerate}

\textbf{练习 B：文学评论句式纠错与 C1 凝练升级（修正语法与搭配错误，或按要求改写）}
\begin{enumerate}
  \item 【语序纠错】\textit{Ved nærmere eftersyn, eventyret er ikke blot en sød fortælling for børn.} $\rightarrow$
  \item 【搭配纠错】\textit{Litteraturforskere fortolker ofte værket til en skarp allegori af det borgerlige samfund.} $\rightarrow$
  \item 【复数词形纠错】\textit{H.C. Andersen skrev mange berømte eventyrer, som stadig diskuteres i dag.} $\rightarrow$
  \item 【表里对比升级】用 \textit{Ved første øjekast... Ved nærmere eftersyn...} 改写：\textit{Historien virker først meget enkel, men senere opdager man, at den kritiserer samfundet.} $\rightarrow$
  \item 【否定升级】用 \textit{Snarere end blot at være... læses værket som...} 改写：\textit{Skyggen er ikke kun en uhyggelig historie; den handler i høj grad om tomme facader.} $\rightarrow$
  \item 【潜台词提炼】用 \textit{Under den tilsyneladende... gemmer der sig...} 改写：\textit{Teksten ser harmløs ud, men den indeholder faktisk en bitter kritik af snobberi.} $\rightarrow$
\end{enumerate}

\textbf{练习 C：中译丹（运用文学批评词汇与标准句式翻译下列评论句）}
\begin{enumerate}
  \item 在看似天真的表层之下，这些童话隐藏着尖锐的潜台词。
  \item 经仔细审视，该文本实际上是一出关于虚荣心的辛辣阶级讽刺。
  \item 这部作品不应只被读作儿童故事，而应被视为一则现代寓言。
  \item 安徒生的文本暴露了人们在面对权力时的集体自欺。
  \item 他的许多故事都折射出十九世纪僵化的社会等级。
  \item 只有能够看穿双关多义的读者，才能领会作者的艺术深度。
\end{enumerate}

% ---------------------------------------------------------------------
\ncdsubsec{六、PD3 / Studieprøven 实战微任务与参考答案}

\begin{emailbox}{Studieprøven / PD3 文学文化评论写作任务（220--260 ord）}
\textbf{Opgave}: Mange udlændinge kender kun H.C. Andersens eventyr fra barndommen eller fra romantiske filmudgaver. Skriv en \textbf{sammenhængende litterær og kulturel kommentar på 220--260 ord}, hvor du forklarer, hvorfor hans værker i Danmark ofte læses med en helt anden alvor. Inddrag eksempler fra mindst to af de nævnte eventyr (\textit{Kejserens nye klæder}, \textit{Skyggen} eller \textit{Den grimme ælling}), og brug relevante litterære analysebegreber samt kontraststrukturer fra lektionen.
\end{emailbox}

\begin{answerbox}
\textbf{Før du læser 参考回答}
\begin{enumerate}
  \item \textit{Litteraturforskere peger på, at Andersens eventyr har en bevidst dobbeltbund, der både henvender sig til børn og voksne. ”Kejserens nye klæder” udstiller konformitet, smiger og et kollektivt selvbedrag ved hoffet, mens ”Skyggen” fungerer som en mørk allegori over overfladiskhed og opportunisme, der fortrænger integritet.}
  \item \textit{Henrik finder eventyret næsten ubarmhjertigt, mens Sofie påpeger en klar klassesatire i fortællingen, fordi ællingen først respekteres, når omgivelserne ser, at den er en svane. Samtidig mener hun, at man bør læse det som både selvforsvar og kritik af tidens snobberi.}
\end{enumerate}

\textbf{Tekst B 阅读核对答案}
\begin{enumerate}
  \item \da{Forkert}（两人均认为现代重读时会感到某种冷酷，索菲娅更指出其中包含阶级讽刺与多重潜台词）
  \item \da{Først når de andre opdager, at den tilhører en finere art (er en svane).}
  \item \da{Både som et selvforsvar fra en fattig skomagersøn og som en kritik af snobberiet i datidens København.}
\end{enumerate}

\textbf{练习 A 答案}
\begin{enumerate}
  \item \da{dobbeltbund}（文学分析核心词：文本双层意涵，配不定冠词 \textit{en}）
  \item \da{klassesatire}（阶级讽刺作品）
  \item \da{udstiller}（动词：暴露、揭露虚伪或缺点）
  \item \da{som}（固定搭配：\textit{fortolke noget som noget} 将……解读为……）
  \item \da{afspejler}（动词：折射、反映出某种社会现实）
  \item \da{end}（固定升级结构：\textit{Snarere end blot at...} 与其仅仅……不如……）
  \item \da{undertekst}（潜台词、言外之意）
  \item \da{gennemskuer}（动词：看穿、识破反讽或表象）
\end{enumerate}

\textbf{练习 B 答案（参考）}
\begin{enumerate}
  \item \textit{Ved nærmere eftersyn er eventyret ikke blot en sød fortælling for børn.}（去掉英语式逗号，且句首状语触发主句 V2 倒装：\textit{er eventyret}）
  \item \textit{Litteraturforskere fortolker ofte værket som en skarp allegori over det borgerlige samfund.}（搭配修正：\textit{fortolke som} 与 \textit{allegori over / på}）
  \item \textit{H.C. Andersen skrev mange berømte eventyr, som stadig diskuteres i dag.}（中性名词 \textit{et eventyr} 单复同形，\textit{eventyrer} 意为“冒险家”）
  \item \textit{Ved første øjekast virker historien meget enkel, men ved nærmere eftersyn kritiserer den samfundet.}（运用“乍看……细看……”标准对照结构提升思辨层次）
  \item \textit{Snarere end blot at være en uhyggelig historie læses Skyggen som en allegori over tomme facader.}（用 \textit{Snarere end blot at...} 实现表里定性升级）
  \item \textit{Under den tilsyneladende harmløse form gemmer der sig en bitter kritik af snobberi.}（用形式主语存现结构突出潜台词焦点）
\end{enumerate}

\textbf{练习 C 答案}
\begin{enumerate}
  \item \textit{Under den tilsyneladende uskyldige overflade gemmer disse eventyr en skarp undertekst.}
  \item \textit{Ved nærmere eftersyn er teksten faktisk en bidende klassesatire over forfængelighed.}
  \item \textit{Dette værk bør ikke kun læses som børnelitteratur, men som en moderne allegori.}
  \item \textit{Andersens tekst udstiller menneskers kollektive selvbedrag i mødet med magten.}
  \item \textit{Mange af hans historier afspejler de stive sociale skel i 1800-tallet.}
  \item \textit{Kun den læser, der kan gennemskue tvetydigheden, forstår forfatterens kunstneriske dybde.}
\end{enumerate}

\textbf{PD3 / Studieprøven 高分范文（约 250 词）}\\
\textit{\textbf{Mellem barnets blik og voksensatire: H.C. Andersens skjulte verden}}

\textit{For mange læsere uden for Danmark står H.C. Andersen først og fremmest som en mild eventyrdigter, hvis historier hører hjemme i barndommens trygge rammer. \dabold{Ved nærmere eftersyn rummer} hans forfatterskab imidlertid en bevidst \da{dobbeltbund}, der retter sig direkte mod voksne og kritiske læsere. Under den tilsyneladende enkle og legende form gemmer der sig nemlig en skarp samfundskritik og en filosofisk \da{tvetydighed}, som gør teksterne vedvarende aktuelle også i det moderne samfund.}

\textit{Dette dobbelte blik træder særligt tydeligt frem i \textit{Kejserens nye klæder}. \dabold{Ved første øjekast} virker fortællingen som en munter anekdote om en forfængelig monark, men i virkeligheden fungerer den som en præcis \da{klassesatire}. Ved at beskrive et hof, hvor ingen tør indrømme det åbenlyse af frygt for at tabe ansigt eller stilling, \da{udstiller} Andersen et kollektivt \da{selvbedrag} skabt af konformitet og magtangst. Det er ikke blot kejserens \da{forfængelighed}, der afsløres, men hele elitens mangel på moralsk mod.}

\textit{Endnu mere foruroligende er \textit{Skyggen}. \dabold{Snarere end blot at være} en dyster gyserhistorie læses værket i dag som en tidløs \da{allegori} over, hvordan ydre overflade og kynisk imagepleje kan fortrænge sandhed og integritet. Hvor skyggen ubesværet stiger i graderne gennem tom charme, går den lærde til grunde med sine idealer intakte. På tilsvarende vis \da{afspejler} selv \textit{Den grimme ælling} en dyb og smertelig erfaring med social udstødelse og snobberi.}

\textit{Netop derfor bør Andersens værker aldrig reduceres til ufarlig nostalgi. Først når man \da{gennemskuer} den \da{flertydige} \da{undertekst}, opdager man, hvor ubarmhjertig --- og samtidig befriende ærlig --- hans litterære kunst egentlig er.}

\small\textbf{名师点评：}这篇约 250 词的范文严格契合 C1 文学文化评论写作范式：开篇点出跨文化印象与丹麦本土学术视角之别；中段紧扣《皇帝的新衣》《影子》与《丑小鸭》三篇经典，自然调度 \textit{dobbeltbund, klassesatire, selvbedrag, allegori, undertekst} 等核心批评词汇；全篇句型节奏张弛有度，精准呈现表里对照（\textit{Ved første øjekast / Ved nærmere eftersyn / Snarere end blot at...}）与主句 V2 倒装，结尾升华点题，毫无生硬拼凑感。
\end{answerbox}
